\documentclass[12pt,a4paper]{article}

\usepackage[utf8]{inputenc}
\usepackage[T1]{fontenc}
\usepackage[brazil]{babel}
\usepackage[top=3cm,left=3cm,bottom=2cm,right=2cm]{geometry}
\usepackage{mathptmx}
\usepackage{inconsolata}
\usepackage{setspace}
\usepackage{indentfirst}
\usepackage{titlesec}
\usepackage{booktabs}
\usepackage{array}
\usepackage{listings}
\usepackage{xcolor}
\usepackage{xurl}
\urlstyle{same}

\onehalfspacing
\setlength{\parindent}{1.25cm}

\titleformat*{\section}{\normalfont\bfseries\normalsize}
\titleformat*{\subsection}{\normalfont\bfseries\normalsize}

\lstset{
    basicstyle=\ttfamily\footnotesize,
    breaklines=true,
    columns=fullflexible,
    keepspaces=true,
    frame=single,
    rulecolor=\color{black!25},
    backgroundcolor=\color{black!4},
    xleftmargin=0.2cm,
    xrightmargin=0.2cm,
    framexleftmargin=4pt,
    framexrightmargin=4pt,
    framextopmargin=4pt,
    framexbottommargin=4pt,
    literate=
        {á}{{\'a}}1 {é}{{\'e}}1 {í}{{\'i}}1 {ó}{{\'o}}1 {ú}{{\'u}}1
        {Á}{{\'A}}1 {É}{{\'E}}1 {Í}{{\'I}}1 {Ó}{{\'O}}1 {Ú}{{\'U}}1
        {ã}{{\~a}}1 {õ}{{\~o}}1 {Ã}{{\~A}}1 {Õ}{{\~O}}1
        {â}{{\^a}}1 {ê}{{\^e}}1 {ô}{{\^o}}1
        {à}{{\`a}}1 {ç}{{\c{c}}}1 {Ç}{{\c{C}}}1
}

\begin{document}

\section{Questão 12: Agente com LLM para planejamento acadêmico}

\subsection{Aplicação e interface}

O assistente monta a agenda de um professor para um dia. O usuário descreve em português a janela disponível, as tarefas, suas durações, prioridades e compromissos fixos. O agente pede esclarecimentos quando faltam dados. Antes de calcular, mostra a interpretação para que o usuário possa corrigir os dados ou confirmar o cálculo.

\begin{table}[htbp]
    \centering
    \begin{tabular}{>{\raggedright\arraybackslash}p{4.5cm} >{\raggedright\arraybackslash}p{10.3cm}}
        \toprule
        \textbf{Componente} & \textbf{Função} \\
        \midrule
        Interface em linguagem natural & Recebe o pedido e seus complementos no terminal. Exibe as tarefas interpretadas, o plano e as pendências. \\
        \addlinespace
        Inferência do LLM & O modelo local, servido pelo Ollama, converte o pedido em JSON conforme um esquema. Extrai janela, duração, prioridade, prazo e início fixo [1]. \\
        \addlinespace
        Planejador em Python & Confere valores e conflitos, separa a janela em blocos livres e calcula a agenda de maior utilidade. \\
        \addlinespace
        Explicabilidade & Mostra a utilidade total, os horários, a prioridade de cada tarefa selecionada e o motivo das pendências. \\
        \bottomrule
    \end{tabular}
\end{table}

\subsection{Planejamento e critérios}

Cada compromisso fixo é obrigatório. O cálculo é interrompido se houver sobreposição, horário fora da janela ou compromisso fixo que termine após seu limite. As tarefas flexíveis só entram no plano se couberem inteiras em um bloco e terminarem até a hora-limite, se houver.

O planejador compara as sequências possíveis dentro de cada bloco. Em cada passo pode colocar uma tarefa no início do espaço restante ou passar ao próximo bloco. Guarda o melhor resultado para estados já visitados, definidos pelo bloco, pelo horário e pelas tarefas restantes. Como as tarefas flexíveis podem começar a qualquer momento em seu bloco, compactar cada sequência à esquerda preserva a viabilidade.

A utilidade é a soma dos pesos das tarefas concluídas. As prioridades alta, média e baixa valem 10, 5 e 1. Em empate, o agente favorece o plano que conclui mais tarefas com prazo para hoje. O prazo de amanhã aparece na pendência para orientar o próximo planejamento.

O usuário pode indicar uma hora-limite para uma tarefa no dia planejado. Por exemplo, se um relatório de 60 minutos vence às 15h e a única janela livre começa às 15h, o sistema deixa o relatório pendente e informa que nenhum bloco o comporta até o horário solicitado.

\subsection{Demonstração do cenário da questão 11}

\noindent Pedido de exemplo

``Tenho das 14h às 18h. Preciso preparar uma aula de IA hoje em 90 minutos, prioridade alta, corrigir trabalhos em 120 minutos até amanhã, prioridade média, e participar de uma reunião de pesquisa das 16h às 17h, prioridade alta. Também tenho um artigo para revisar em 60 minutos até a próxima semana, prioridade baixa.''

\begin{table}[htbp]
    \centering
    \begin{tabular}{>{\raggedright\arraybackslash}p{3.6cm} >{\raggedright\arraybackslash}p{9.5cm}}
        \toprule
        \textbf{Horário} & \textbf{Resultado} \\
        \midrule
        14h00--15h30 & Preparar aula de IA (10 pontos). \\
        \addlinespace
        15h30--16h00 & Tempo livre na agenda. \\
        \addlinespace
        16h00--17h00 & Reunião de pesquisa (10 pontos), horário fixo. \\
        \addlinespace
        17h00--18h00 & Revisar artigo (1 ponto). \\
        \addlinespace
        Pendente & Corrigir trabalhos. Prazo até amanhã. \\
        \bottomrule
    \end{tabular}
\end{table}

A correção exige duas horas seguidas. Só há um bloco desse tamanho, antes da reunião, e ele também é o único que comporta a aula. Ao usar esse bloco para a aula, a utilidade do plano passa de 16 para 21 pontos. Os 30 minutos livres antes da reunião são insuficientes para a correção. A justificativa combina os dados de duração, horários fixos e pesos do planejamento.

\subsection{Execução e limites}

Em um computador com Python e Ollama disponíveis, baixe o modelo com \texttt{ollama pull qwen2.5:3b}. Com o serviço local em execução, rode \texttt{python agenda.py} na pasta da questão e digite o pedido. O argumento \texttt{-{}-modelo} permite selecionar outro modelo já instalado.

O comando \texttt{python agenda.py -{}-demo} reproduz o resultado da tabela com dados estruturados de exemplo. A execução do modelo exige um servidor Ollama ativo. O escopo cobre um dia, até oito tarefas e um recurso, o tempo do professor. A confirmação da interpretação permite identificar erros de extração antes de montar a agenda.

\section*{Referências}

\noindent [1] OLLAMA. \textbf{Structured Outputs}. Documentação oficial. Disponível em: \url{https://docs.ollama.com/capabilities/structured-outputs}.

\end{document}
